\subsection{Objective and Input Image}

Task 1 requires displaying the original chip image, applying contrast
enhancement, and commenting on the result.  The input image is an RGB bitmap of
size $990\times367$ pixels.  Its RGB pixel range is 41--255, while the grayscale
luminance range is approximately 41.91--252.31.  The grayscale median is 82.91,
and the 2nd and 98th percentiles are 63.63 and 156.00 respectively.

The original image in \cref{fig:task1-comparison} shows bright characters on a
dark, non-uniform background.  The main difficulty is not low brightness alone:
the background contains dense ring-like texture and bright speckle noise.  A
useful enhancement method must therefore improve the character-to-background
separation without making the background texture look like character strokes.

\subsection{Implemented Enhancement Methods}

The enhancement operators follow standard digital image-processing
definitions~\cite{gonzalez2018}; their experiment code was implemented in Python
in \texttt{task\_1/task\_1.py}.  The
RGB image was first converted to grayscale luminance by
\begin{equation}
  Y = 0.299R + 0.587G + 0.114B .
  \label{eq:task1-grayscale}
\end{equation}
The grayscale image is the baseline for comparison.  Three contrast enhancement
methods were tested on the same baseline:

\begin{enumerate}
  \item \textbf{2--98\% percentile stretching.}  The 2nd percentile is mapped to
  0 and the 98th percentile is mapped to 255:
  \begin{equation}
    I_{\mathrm{out}} =
    \mathrm{clip}\!\left(
      255\frac{I-p_{2}}{p_{98}-p_{2}},\,0,\,255
    \right).
    \label{eq:task1-percentile}
  \end{equation}
  This ignores extreme dark and bright outliers and uses the available grayscale
  range more aggressively.

  \item \textbf{Gamma correction.}  A power-law transform with $\gamma=1.5$ was
  applied:
  \begin{equation}
    I_{\mathrm{out}} = 255\left(\frac{I}{255}\right)^{1.5}.
    \label{eq:task1-gamma}
  \end{equation}
  Since $\gamma>1$, mid-level background pixels are darkened.  This tests whether
  suppressing the background improves character visibility.

  \item \textbf{Global histogram equalization.}  A 256-level cumulative
  distribution function mapping was applied to spread the global grayscale
  histogram.  This is included as a strong global enhancement baseline.
\end{enumerate}

For a fair visual comparison, all grayscale images in
\cref{fig:task1-comparison} are displayed with the same scale,
$v_{\min}=0$ and $v_{\max}=255$.  No method is given a separate automatic display
stretch.

\begin{figure}[!htbp]
  \centering
  \includegraphics[width=\textwidth]{task1_comparison.png}
  \caption{Original image and contrast enhancement results.  All grayscale
  images use the same display range of 0--255.}
  \label{fig:task1-comparison}
\end{figure}

\subsection{Evaluation Protocol}

The result was evaluated both visually and quantitatively.  Visual evaluation
checks whether the characters, especially \texttt{T12} and
\texttt{HD44780A00}, become clearer without excessive background amplification.
Quantitative evaluation uses three values:

\begin{itemize}
  \item \textbf{CNR}: contrast-to-noise ratio between fixed character-stroke
  patches and fixed nearby background patches,
  \begin{equation}
    \mathrm{CNR} =
    \frac{|\mu_c-\mu_b|}
    {\sqrt{\sigma_c^2+\sigma_b^2}} .
    \label{eq:task1-cnr}
  \end{equation}
  Larger CNR means better separation between character pixels and background
  pixels in the selected regions.

  \item \textbf{Otsu separability}: normalized between-class variance from
  Otsu's thresholding criterion.  A larger value suggests that the grayscale
  histogram is easier to split into two classes, but it does not guarantee that
  the foreground class corresponds only to characters.

  \item \textbf{Saturation ratio}: percentage of pixels clipped to the endpoints
  0 or 255.  A large value indicates loss of grayscale detail.
\end{itemize}

\begin{table}[!htbp]
  \centering
  \caption{Task 1 quantitative comparison.}
  \label{tab:task1-metrics}
  \begin{tabular}{lccc}
    \toprule
    Method & CNR & Otsu separability & Saturation ratio \\
    \midrule
    Grayscale baseline & 0.6933 & 0.6941 & 0.00\% \\
    2--98\% percentile stretch & \textbf{0.7013} & 0.7181 & 4.22\% \\
    Gamma correction, $\gamma=1.5$ & 0.6948 & 0.7077 & 0.00\% \\
    Global histogram equalization & 0.5669 & \textbf{0.7410} & 0.41\% \\
    \bottomrule
  \end{tabular}
\end{table}

\begin{figure}[!htbp]
  \centering
  \includegraphics[width=0.95\textwidth]{task1_histograms.png}
  \caption{Histograms of the grayscale baseline and the three enhancement
  methods.}
  \label{fig:task1-histograms}
\end{figure}

\subsection{Result Analysis}

The grayscale baseline preserves the original information, but the characters
are only weakly separated from the textured background.  Its CNR is 0.6933 and
its Otsu separability is 0.6941, which provides the reference level for the
enhancement methods.

The 2--98\% percentile stretch gives the most useful result for this task.  It
raises CNR from 0.6933 to 0.7013 and raises Otsu separability from 0.6941 to
0.7181.  In \cref{fig:task1-comparison}, the strokes in \texttt{T12} and
\texttt{HD44780A00} become visibly brighter and easier to locate.  The trade-off
is a 4.22\% saturation ratio, caused by clipping the lowest and highest tails of
the intensity distribution.  This clipping is acceptable here because it mainly
increases character visibility; it does not dominate the image with false
foreground structures.

Gamma correction with $\gamma=1.5$ is conservative.  It slightly improves CNR
to 0.6948 and Otsu separability to 0.7077 without endpoint clipping.  However,
the improvement is small because the transform darkens both background texture
and weaker character strokes.  It does not separate the characters as clearly as
percentile stretching.

Global histogram equalization is not selected despite its highest Otsu
separability of 0.7410.  The visual result shows that it strongly amplifies the
ring-like background texture and bright speckles.  Its CNR also drops to 0.5669,
which means the selected character and background patches become less separable
than in the original grayscale image.  This is a clear example where a wider or
more separable global histogram does not imply better character readability.
For later binary conversion and segmentation, this result would likely create
many false foreground regions.

\subsection{Selected Output for Later Tasks}

The selected Task 1 output is the 2--98\% percentile-stretched image saved in
\texttt{task\_1/outputs/}.  It provides the best balance among the tested
methods: character strokes become clearer, the measured CNR is the highest, and
the background is not amplified as severely as in global histogram equalization.
This output is therefore the preferred input candidate for the following
smoothing, thresholding, outline detection, and segmentation stages.

\FloatBarrier
